\chapter*{Abstract}
\addcontentsline{toc}{chapter}{Abstract}

Sepsis is a leading cause of intensive-care mortality, and earlier recognition
improves survival. Research on sepsis prediction is dominated by machine
learning and, on the MIMIC-IV database specifically, by prognostic modelling of
mortality \emph{after} diagnosis rather than prediction of \emph{onset} itself.
This thesis develops and validates a fully interpretable, purely statistical
pipeline for real-time Sepsis-3 onset prediction. Using 65{,}241 adult ICU
stays from MIMIC-IV v3.1, each stay was resampled to an hourly grid of vitals
and laboratory values, hourly Sequential Organ Failure Assessment (SOFA) scores
were computed, and the exact onset hour was labelled under the Sepsis-3
consensus definition. Continuously updated risk was obtained through a dynamic
landmark supermodel rather than any neural or ensemble method: a single
elastic-net logistic model fitted on stacked landmark datasets across six
prediction hours, achieving a pooled out-of-sample AUROC of 0.797.

Three methodological commitments distinguish the work. First, the
proportional-hazards assumption underlying Cox regression was formally tested
using Schoenfeld residuals and, because the global test failed, a time-varying
reformulation was fitted in response. Second, incident (new-onset) sepsis was
separated from prevalent sepsis using an hour-six landmark, which revealed that
standard bedside scores range from non-discriminative to inverse on this task
(SOFA area under the receiver operating characteristic curve 0.36, qSOFA 0.49,
SIRS 0.59), whereas a multivariable landmark nomogram attained 0.775. Third, an
informative-missingness analysis showed that whether a laboratory test was
ordered is a stronger predictor than its measured value, quantifying a
label-recognition ceiling intrinsic to retrospectively labelled ICU data.
Beyond this internal analysis, the frozen model was externally validated on the
eICU Collaborative Research Database of more than 200 United States ICUs: discrimination
transported with modest degradation (area under the curve 0.686 against an
internal eICU ceiling of 0.711), and a single-parameter recalibration corrected
the expected miscalibration. All extraction, modelling, and evaluation code is
released for reproducibility.
